\chapter{Architecture globale}
\label{ch:archi}

\section{Vue logique}

L'architecture suit un modèle en couches classique, enrichi d'une couche d'observabilité. La figure~\ref{fig:archi-logique} synthétise les composants et leurs interactions.

\begin{figure}[htbp]
  \centering
  \begin{tikzpicture}[node distance=0.9cm and 1.2cm]
    \node[block, cloud, minimum width=3.2cm] (nav) {Navigateur\\(utilisateur)};
    \node[block, svc, below=of nav, minimum width=3.6cm] (web) {Frontend React\\nginx port 80};
    \node[block, svc, below=of web, minimum width=3.6cm] (api) {API FastAPI\\port 8000};
    \node[block, data, below=of api, minimum width=3cm] (pg) {PostgreSQL\\16};
    \node[block, infra, right=2.8cm of api, minimum width=3.2cm] (prom) {Prometheus\\Grafana};

    \draw[line] (nav) -- (web);
    \draw[line] (web) -- node[right, font=\footnotesize] {REST /api/v1} (api);
    \draw[line] (api) -- (pg);
    \draw[dashedline] (api) -- node[above, font=\footnotesize] {/metrics} (prom);
  \end{tikzpicture}
  \caption{Architecture logique de StockFlow}
  \label{fig:archi-logique}
\end{figure}

\section{Deux topologies de déploiement}

\subsection{Mode développement (Docker Compose)}

Le fichier \texttt{docker/compose.dev.yml} définit un réseau bridge \texttt{pfa} reliant postgres, api, web et \textbf{Adminer} (visualisation BDD, port~8080). Ports publiés : 3000 (UI), 8000 (API), 5432 (PostgreSQL), 8080 (Adminer).

\subsection{Mode validation Kubernetes (minikube)}

Le cluster \textbf{minikube} (profil \texttt{pfa}) est provisionné par \texttt{make cluster} : addon ingress, kube-prometheus-stack, pgAdmin, dashboard Kubernetes, chart Helm \texttt{pfa-stock} (\texttt{values-minikube.yaml}). Les hôtes \texttt{pfa.test}, \texttt{api.pfa.test}, \texttt{grafana.test} et \texttt{pgadmin.test} sont mappés via \texttt{/etc/hosts} sur l'IP minikube.

\subsection{Production optionnelle (VPS k3s)}

Un déploiement sur VPS reste documenté (\texttt{docs/vps-setup.md}, workflow \texttt{deploy.yml} manuel) mais n'est pas requis pour la validation du PFA.

\begin{figure}[htbp]
  \centering
  \begin{tikzpicture}[font=\small\sffamily, node distance=0.9cm and 1.4cm]
    \node[block, infra, minimum width=3.4cm] (ing) {ingress-nginx\\TLS 443};
    \node[block, svc, below=of ing, minimum width=4.2cm] (app) {namespace pfa-stock\\web, api, postgres};
    \node[block, infra, right=of ing, minimum width=3.2cm] (mon) {monitoring\\Grafana, Prometheus};
    \draw[line] (ing) -- (app);
    \draw[dashedline] (app.east) -- node[above, font=\footnotesize] {métriques} (mon.west);
  \end{tikzpicture}
  \caption{Déploiement minikube (vue simplifiée)}
  \label{fig:k8s}
\end{figure}

Les scripts \texttt{demarrer-cluster-minikube.sh} installent la plateforme ; le job CI \texttt{kubernetes} reproduit le déploiement Helm sur un minikube éphémère.

\section{Flux de données principaux}

\subsection{Authentification}

\begin{enumerate}
  \item Le client envoie \texttt{POST /api/v1/auth/connexion} avec email et mot de passe.
  \item L'API vérifie le hash bcrypt, émet un JWT.
  \item Le frontend stocke le jeton dans \texttt{sessionStorage}.
\end{enumerate}

\subsection{Mouvement de stock}

Le service \texttt{gestion\_stock.py} encapsule la logique transactionnelle (chapitre~\ref{ch:backend}).

\subsection{Observabilité}

L'instrumentateur Prometheus expose \texttt{/metrics}. Sur minikube, Grafana est accessible sur \texttt{http://grafana.test} (après entrée \texttt{/etc/hosts}).

\section{Chaîne d'outillage DevOps}

\begin{figure}[htbp]
  \centering
  \begin{tikzpicture}[node distance=0.55cm]
    \node[block, minimum width=2.8cm] (git) {Git / GitHub};
    \node[block, minimum width=2.8cm, right=1.1cm of git] (ci) {GitHub Actions\\CI + CD};
    \node[block, minimum width=2.8cm, right=1.1cm of ci] (ghcr) {GHCR};
    \node[block, minimum width=2.8cm, below=1cm of ghcr] (helm) {Helm};
    \node[block, minimum width=2.8cm, left=1.1cm of helm] (mk) {minikube};

    \draw[line] (git) -- (ci);
    \draw[line] (ci) -- (ghcr);
    \draw[line] (ghcr) -- (helm);
    \draw[line] (helm) -- (mk);
  \end{tikzpicture}
  \caption{Chaîne d'outillage (validation CI + démo minikube)}
  \label{fig:devops-chain}
\end{figure}

\section{Choix architecturaux structurants}

\begin{table}[htbp]
  \centering
  \caption{Décisions architecturales et justification}
  \label{tab:decisions-archi}
  \begin{tabularx}{\textwidth}{@{}lX@{}}
    \toprule
    \textbf{Décision} & \textbf{Justification} \\
    \midrule
    API stateless & Scaling horizontal des pods API ; session portée par JWT. \\
    Frontend statique (nginx) & Surface d'attaque réduite en cluster. \\
    PostgreSQL relationnel & Transactions ACID pour les mouvements. \\
    minikube + Helm & Validation K8s locale et en CI sans VPS dédié. \\
    GHCR + Helm (optionnel) & Images versionnées ; déploiement VPS manuel. \\
    kube-prometheus-stack & Observabilité (Grafana, ServiceMonitor). \\
    \bottomrule
  \end{tabularx}
\end{table}
